\documentclass[11pt]{article}
\usepackage{fontspec}
\usepackage[french]{babel}
\usepackage{amsmath}
\usepackage{amssymb}
\usepackage{amsthm}
\usepackage{mathrsfs}
\usepackage{tikz-cd}
\usepackage{enumitem}

\newcommand{\printedpage}[1]{\marginpar{\footnotesize p.~#1}}
\newcommand{\msp}[1]{\marginpar{\footnotesize [#1]}}
\newcommand{\Fl}{\mathrm{Fl}}
\newcommand{\Is}{\mathrm{Is}}
\newcommand{\Fib}{\mathrm{Fib}}

\begin{document}

\printedpage{54}
% head: A. Grothendieck : Dérivateurs, Ch. II, Édition des Universités de Montpellier II et Paris VII
On est bien sûr intéressé surtout aux éléments clos de $\mathcal{P}$, $\mathcal{Q}$, i.e. les parties $A$, $L$ de $\Fl(\mathcal{M})$, contenant les isomorphismes, stables par composition (et $L$ stable aussi par changement de base, $A \subseteq W$) telles que l'on ait
\[
A = \mathbf{A}_W(L), \qquad L = \mathbf{L}(A) \tag*{(6.8)}
\]
\texttt{[plutôt }$\mathbf{L}_W(A)$\texttt{]}.

\msp{page 87}

Je garderai les notations $\natural$, $\flat$ du paragraphe précédent. Quels sont les deux couples extrêmes ?

Partant de $A_0 = \Is(\mathcal{M})$ (plus petit des éléments de $\mathcal{P}$), on trouve $\mathbf{L}_W(A_0) = \Fl(\mathcal{M})$ (plus grand élément de $\mathcal{Q}$), et $A_0^{\natural} = \mathbf{A}_W(\Fl(\mathcal{M})) = W_{\mathrm{univ}}$, ensemble des flèches qui sont dans $W$ et le restent par tout changement de base :
\[
A_0^{\natural} = W_{\mathrm{univ}}, \qquad \beta_0(A_0) \overset{\text{déf}}{=} \mathbf{L}_W(A_0) = \Fl(\mathcal{M}), \qquad \text{\scriptsize où $A_0 = \Is(\mathcal{M})$ plus petit élément de $\mathcal{P}$.} \tag*{(6.9)}
\]
Inversement, partons de $L_0 = \Is(\mathcal{M})$, le plus petit élément de $\mathcal{Q}$. On a $\mathbf{A}_W(L_0) = W$ ($W$ étant stable par composition et contenant les isomorphismes), et $L_0^{\flat} = \mathbf{L}_W(W) = \Fib_W$, les $W$-fibrations :
\[
L_0^{\flat} = \Fib_W, \qquad \alpha_0(L_0) = W, \qquad \text{\scriptsize où $L_0 = \Is(\mathcal{M})$, plus petit élément de $\mathcal{Q}$.} \tag*{(6.10)}
\]
Donc les parties $\mathcal{R}$-closes à gauche de $\Fl(\mathcal{M})$ se promènent entre $W_{\mathrm{univ}}$ et $W$, les parties $\mathcal{R}$-closes à droite entre $\Fib_W$ et $\Fl(\mathcal{M})$.

\textbf{Proposition 6.1.} \textit{Soient $A \in \mathcal{P}$, $L \in \mathcal{Q}$ tels que $A = \alpha_0(L)$, $L = \beta_0(A)$ (i.e. $A = \mathbf{A}_W(L)$, $L = \mathbf{L}_W(A)$). Pour que $L$ soit stable par composition, il faut et il suffit que $A$ soit stable par changement de base par les flèches $f \in L$. (Plus généralement, si $L \in \mathcal{Q}$ est un élément quelconque de $\mathcal{Q}$, et si on suppose de} \msp{page 88} \textit{plus que $L$ est stable par composition, alors $A = \mathbf{A}_W(L)$ est stable par changement de base par $L$.)}

Terminologie : pour $A \in \mathcal{P}$, i.e. $A$ une partie de $\Fl(\mathcal{M})$ contenue dans $W$ contenant les isomorphismes et stable par composition, $\mathbf{L}_W(A)$ s'appelle le \textit{transporteur de $A$ dans $W$}. (Il y a une notion duale de \textit{cotransporteur.}) C'est un élement de $\mathcal{Q}$, donc une partie de $\Fl(\mathcal{M})$ contenant les isomorphismes et stable par changement de base. D'autre part, pour $L$ dans $\mathcal{Q}$, les $f \in \mathbf{A}_W(L)$ s'appellent flèches \textit{$W$-$L$-asphériques,} ou les \textit{$W$-$L$-équivalences,} ou les \textit{$W$-équivalences $L$-universelles,} et $\mathbf{A}_W(L)$ s'appelle le \textit{$L$-noyau de $W$}. On pourrait le noter plus simplement $W_L$ (quand cette notation ne conflicte pas avec $W_S \subseteq \Fl(\mathcal{M}/S)$ pour $S \in \mathrm{Ob}\,\mathcal{M}$). C'est une partie de $W$ qui (comme $W$ lui-même) contient les isomorphismes et est stable par composition. % [sic] élement -> élément, missing accent

Enfin, il y a lieu aussi d'introduire le \textit{stabiliseur} $\mathrm{Stab}(A)$, ou transporteur de $A$ dans $A$, qui ne dépend pas de la donné de $W$, mais seulement de la partie $A$ de $\Fl(\mathcal{M})$, contenant les isomorphismes et stable par composition. On pourrait aussi le noter $\Fib_A$, et appeler ses éléments des \textit{$A$-fibrations} ($A$ jouant à présent le rôle de $W$). On a $\mathrm{Stab}(A) \subseteq \mathbf{L}_W(A)$, et il semblerait qu'en % [sic] donné -> donnée, feminine noun
% foot: Version 8 juillet 2001
\printedpage{55}
% head: A. Grothendieck : Dérivateurs, Ch. II, Édition des Universités de Montpellier II et Paris VII
\msp{page 89}
pratique, ce sont ces stabiliseurs, plus que les transporteurs $\mathbf{L}_W(A)$, qui sont importants (quand les deux ne coïncident pas). C'est ainsi que dans Cat, si $W$ est un localiseur fondamental, les morphismes $W$-lisses (resp. $W$-propres) $\mathrm{Liss}_W$ ($\mathrm{Prop}_W$), qui sont les stabiliseurs des ensembles $W_{\mathrm{as}}$ ($W_{\mathrm{coas}}$) des ensembles de morphismes $W$-asphériques (resp. $W$-coasphériques), jouant un rôle essentiel, alors que je ne sais pas caractériser de façon sensible les flèches des transporteurs $\mathbf{L}_W(W_{\mathrm{as}})$, $\mathbf{L}_W(W_{\mathrm{coas}})$, lesquels contiennent, en plus des morphismes $W$-lisses (resp. $W$-propres) les $W$-fibrations $\Fib_W$. Je ne sais pas comment caractériser une classe de morphismes qui apparaît aussi hétéroclite! De même, il y a lieu d'introduire et de caractériser le stabiliseur de $W_{\mathrm{bias}}$ (ensemble des flèches $W$-biasphériques), comme les morphismes `$W$-nets' $\mathrm{Net}_W$, par une description similaire à celle des $\mathrm{Liss}_W$, $\mathrm{Prop}_W$, utilisant les bilocalisés $\backslash\underset{u}{X}/$, $\backslash\underset{v}{Y}/$ , au lieu des localisés $X/x$, $Y/y$ ou des colocalisés $x\backslash X$, $y\backslash Y$. Il y a lieu enfin de déterminer le stabiliseur de l'ensemble Equ des Cat-équivalences (ce qui ne dépend pas de la donné d'un $W$ - mais on a $\mathrm{Equ} \subseteq \mathcal{P}(W)$ pour tout $W$ fondamental). % [sic] donné -> donnée, feminine noun
% foot: Version 8 juillet 2001

\printedpage{56}
% head: A. Grothendieck : Dérivateurs, Ch. II, Édition des Universités de Montpellier II et Paris VII
\msp{page 90}

\section*{7\quad Morphismes cofiltrants et localement cofiltrants}

Mais je ne veux pas m'attarder à présent sur ces développements, qui ne me paraissent pas urgents, et voudrais revenir aux questions liées à la notion de catégories filtrantes ou cofiltrantes.

\textbf{Definition 7.1.} Soit $f : X \longrightarrow Y$ une flèche dans Cat. On dit que $f$ est \textit{cofiltrante} (resp. \textit{localement cofiltrante}, resp. \textit{faiblement cofiltrante} (\footnote[64]{J'ai modifié la terminologie : maintenant faiblement cofiltrant $=$ 0-connexe $+$ localement faiblement cofiltrant, \textit{stable par passage} à $X'$ sur $X$, $X'$ étale et $\infty$-connexe. Revoir la terminologie dans le cas faiblement cofiltrant, qui est une notion traîtresse, faute d'être stable par passage \texttt{[à]} un $X'$ étale sur $X$.}), resp. \textit{localement faiblement cofiltrante} ou encore \textit{précofiltrante}) si pour tout $y$ dans $Y$, la catégorie colocalisé $y\backslash X$ est cofiltrante (resp. \dots). (Dualement, on dit que $f$ est \textit{filtrante}, resp. \dots, si les catégories $X/y$ sont filtrantes, resp. \dots, i.e. si $f^o : X^o \longrightarrow Y^o$ est cofiltrante resp. \dots.) % [sic] Definition -> Définition, missing accent; colocalisé -> colocalisée, agreement with catégorie

\textbf{Remarque 7.2.} Notions stables par changement de base coétale. Contient les isomorphismes et même les Cat-équivalences. Les morphismes \textit{étales} sont localement filtrants.

\textbf{Proposition 7.3.} \textit{Soit $f : X \longrightarrow Y$ un morphisme cofiltrant (resp. localement cofiltrant, resp. faiblement cofiltrant, resp. précofiltrant), et supposons que $Y$ soit cofiltrant (resp. localement cofiltrant, resp. faiblement cofiltrant, resp. précofiltrant). Alors $X$ lui-même est cofiltrant (resp. \dots).}

\textsc{Démonstration}. Il s'agit, dans les quatre cas envisagés, de vérifier pour $X$ la condition $\underbrace{\mathrm{PS}_g1' \;/\; \mathrm{PS}_g1}_{\text{deux derniers cas}} \;/\; \mathrm{PS}_g1' + \mathrm{PS}_g2 \;/\; \mathrm{PS}_g1 + \mathrm{PS}_g2$. Mais \msp{page 91} pour chacune des trois conditions $\mathrm{PS}_g1$, $\mathrm{PS}_g1'$, $\mathrm{PS}_g2$ séparément, on voit immédiatement que si elle est vérifiée pour $Y$ et pour les $y\backslash X$, elle l'est pour $X$, d'où la proposition.

\textbf{Corollaire 7.3.} \textit{Les quatre notions envisagées dans (7.1) sont stables par composition.}

\textsc{Démonstration} immédiate, en utilisant la remarque (7.2).
\[
\begin{tikzcd}
X \arrow[d] & z\backslash X \arrow[l] \arrow[d] \\
Y \arrow[d] & z\backslash Y \arrow[l] \arrow[d] \\
Z & z\backslash Z. \arrow[l]
\end{tikzcd}
\]
% foot: Version 8 juillet 2001

\printedpage{57}
% head: A. Grothendieck : Dérivateurs, Ch. II, Édition des Universités de Montpellier II et Paris VII
\textbf{Remarque 7.4.}
\begin{enumerate}
\item[1)] Notons que la notion de morphisme cofiltrant ou faiblement cofiltrant n'est \textit{pas} stable par changement de base quelconque, ni même par une immersion ouverte. Exemple : $e \longrightarrow \boldsymbol{\Delta}^1$, appliquant $e$ en 1, est cofiltrant,
\[
\begin{tikzcd}
 & e \arrow[d, mapsto] \\
0 \arrow[r] & 1
\end{tikzcd}
\]
mais sa fibre en 0 est vide, donc n'est ni cofiltrante, ni faiblement cofiltrante. Je n'ai pas vu d'exemple montrant que la notion de flèche \textit{localement} cofiltrante resp. précofiltrante n'est pas stable par changement de base, et ne voudrais pas m'y attarder, n'ayant pas besoin de connaître la réponse.
\item[2)] Pour que $X$ soit cofiltrante (resp. \dots), il faut et il suffit que $X \longrightarrow e$ le soit. Ainsi la notion absolue est ramenée à une notion relative.
\item[3)] Si $f : X \longrightarrow Y$ est cofiltrante, elle est \msp{page 92} $W$-coasphérique (donc $\in W$) pour tout $W$ satisfaisant l'axiome Loc (8) des limites inductives. En particulier, elle est $W_0$-asphérique, i.e. un foncteur cofiltrant (et même seulement faiblement cofiltrant) est \textit{cofinal.}
\end{enumerate}
La raison principale pour m'intéresser aux flèches cofiltrantes est contenue dans la

\textbf{Proposition 7.4.} \textit{Soit $f : X \longrightarrow Y$ dans} Cat \textit{localement cofiltrant (resp. localement faiblement cofiltrant). Alors :}
\begin{enumerate}
\item[\textit{a)}] \textit{Pour tout $X'$ étale sur $X$, le composé $f' : X' \longrightarrow X \longrightarrow Y$ est localement cofiltrant (resp. \dots), et pour toute factorisation de $f'$ en $X' \xrightarrow{g'} Y' \longrightarrow Y$ avec $Y' \longrightarrow Y$ étale, $g'$ est localement cofiltrant (resp. \dots).}
\item[\textit{b)}] \textit{$\forall\ x$ dans $X$, le morphisme sur les localisés}
\[
f/x : X/x \longrightarrow Y/x \qquad (y = f(x)) % [sic] Y/x -> Y/y, since y = f(x)
\]
\textit{est cofiltrant (resp. faiblement cofiltrant).}
\item[\textit{c)}] \textit{Si $x$, $y$ comme dans b), et $U$ est un ouvert de $X/x$, désignant par $V$ l'ouvert de $Y/y$ engendré par l'image de $f/x$ (i.e. l'ouvert formé des $z$ dans $Y/y$ tels que $z\backslash(X/x)$ soit non vide), la flèche induite $q : U \longrightarrow V$} \texttt{[est]} \textit{cofiltrante (resp. faiblement cofiltrante), a fortiori cofinale (et même $W$-coasphérique, pour tout $W$ satisfaisant l'axiome des limites Loc (8))} (\footnote[65]{Il suffit ici qu'on soit dans l'hypothèse respée ($f$ \textit{faiblement} localement cofiltrant) sauf pour le NB ici \texttt{[`le NB ici' est `(et même ...Loc (8))']}. % [sic] respée -> intended word unclear (perhaps « respectée » or « resp. »), not a French word
}). % ?? uncertain: letter of the induced arrow "q" -- q or g (g is used in g^* below)

\msp{page 93}
\textit{Donc pout tout} cofaisceau \textit{$G$ sur $V$, à valeurs dans une catégorie $\mathcal{M}$ avec $\varinjlim$ , on a} \texttt{[un]} \textit{isomorphisme} % [sic] pout -> pour, misprint
\[
\varinjlim_U g^*(G) \xrightarrow{\ \sim\ } \varinjlim_V G,
\]
% foot: Version 8 juillet 2001
\printedpage{58}
% head: A. Grothendieck : Dérivateurs, Ch. II, Édition des Universités de Montpellier II et Paris VII
\textit{et pour tout préfaisceau $F$ sur $V$, à valeurs dans une catégorie $\mathcal{M}$ avec $\varprojlim$ , on a} \texttt{[un]} \textit{isomorphisme}
\[
\varprojlim_U g^*(F) \xleftarrow{\ \sim\ } \varprojlim_V F.
\]
\end{enumerate}

\textsc{Démonstration}. Je vais d'abord expliciter le

\textbf{Lemme 7.5.} \textit{Soit $h : X' \longrightarrow X$ un morphisme étale dans} Cat. \textit{Alors $h$ est localement cofiltrant. De façon plus précise, pour tout $x$ dans $X$, on a}
\[
\pi_0(x\backslash X') \xrightarrow{\ \sim\ } \mathrm{Ob}\,X'_x,
\]
\textit{toute composante connexe de $x\backslash X'$ contient un élément et un seul de la catégorie discrète $X'_x \subseteq x\backslash X'$. Si $x'$ est dans $X'_x$ (i.e. $x'$ objet de $X'$ tel que $h(x') = x$), c'est un objet} initial \textit{de la composante connexe qui le contient. (Donc cette composante est cofiltrante.)}

Ceci prouve (et au delà) la première assertion de 7.4, compte tenu de 7.3 (stabilité par composition), ou parce qu'une \msp{page 94} catégorie étale (en l'occurrence $y\backslash X'$) au dessus d'une catégorie localement cofiltrante (resp. localement faiblement cofiltrante) (en l'occurrence $y\backslash X$) est itou.
\[
\begin{tikzcd}[row sep=large]
y\backslash X' \arrow[r] \arrow[d, "\text{ét.}"'] \arrow[dr, phantom, "\text{cart.}"] & X' \arrow[d, "\text{ét.}"] \\
y\backslash X \arrow[r] \arrow[d] \arrow[dr, phantom, "\text{cart.}"] & X \arrow[d] \\
y\backslash Y \arrow[r] & Y.
\end{tikzcd}
\]
Pour la deuxième assertion dans a), on peut maintenant supposer $X' = X$. Il faut prouver que pour tout $y'$ dans $Y'$, $y'\backslash X \overset{\text{déf}}{=} X \times_{Y'} y'\backslash Y'$ est localement cofiltrante (resp. localement faiblement cofiltrante). Soit $y$ dans $Y$ l'image de $y'$, et considérons le diagramme de trois carrés cartésiens ci-contre.
\[
\begin{tikzcd}[row sep=large]
y'\backslash X \arrow[r] \arrow[d, "\text{ét.}"] \arrow[dr, phantom, "\textit{cart.}"] & y\backslash X \arrow[r] \arrow[d, "\text{ét.}"] \arrow[dr, phantom, "\textit{cart.}"] & X \arrow[d, "\text{ét.}"] \\
y'\backslash Y' \arrow[r] & y\backslash Y' \arrow[r] \arrow[d] \arrow[dr, phantom, "\textit{cart.}"] & Y' \arrow[d] \\
 & y\backslash Y \arrow[r] & Y
\end{tikzcd}
\]
Le Lemme 7.5 implique que $y'\backslash Y' \longrightarrow y\backslash Y'$ est l'inclusion dans $y\backslash Y'$ d'une composante connexe de $y\backslash Y'$, c'est donc une immersion ouverte. Donc $y'\backslash X \longrightarrow y\backslash X$ est une immersion
% foot: Version 8 juillet 2001

\end{document}
